%% SCRAP: archive/session-logs/2025-11-20-recap
%% SOURCE: docs/working/archive/session-logs/2025-11-20-recap.md
%% STATUS: HISTORICAL
%% FITS: none
%% EDITORIAL: lifted — prose rewritten to press voice

\section{Session Recap: 20 November 2025}
\subsection*{From Factorial DoE to Multivariate Dynamics Analysis}

\textbf{Duration:} approximately 90 minutes.
\textbf{Outcome:} complete paradigm shift documented; ready for code
implementation.
\textbf{Commits:} 3 (design documents plus executive summary).

\subsection{Summary of Paradigm Shift}

Work began with the goal of running the DoE with heartbeat data under a
five-configuration, fifty-runs-per-configuration plan (250 total runs).
Three sequential corrections, each driven by deeper understanding of the
system, produced a fundamentally different experimental framework.

\textbf{Correction~1.} The heartbeat tick interval (\texttt{tick\_ns}) is
not a fixed constant but a tuning variable modulated by the inference engine
in response to load. The experiment therefore measures load-responsive
adaptation, not fixed-interval jitter.

\textbf{Correction~2.} 250 runs are statistically insufficient for
correlation analysis. The revised plan calls for 150 runs per configuration,
for a total of 750 runs (5 configurations $\times$ 150 runs). This provides
adequate statistical power for correlation analysis between metrics.

\textbf{Core insight.} Each metric in the feedback system is simultaneously
a signal and a feedback input. The seven feedback loops form a coupled
dynamical system in which every metric influences the future state of every
other metric. The experimental paradigm therefore shifts from factorial
analysis of main effects to multivariate dynamical systems analysis:
measuring the full trajectory shape of each configuration through
metric-space rather than summarising it by means.

\subsection{Design Documents Produced}

\textbf{MULTIVARIATE\_DYNAMICS\_DESIGN.md (622 lines).}
Defines \texttt{HeartbeatTickSnapshot} (per-tick capture), \\
\texttt{RunStabilityFingerprint} (per-run summary vector), and
\texttt{ConfigurationFingerprint} (per-configuration aggregation). Provides
complete CSV schemas for three output levels, a five-phase analysis pipeline,
and references to control theory and systems-biology literature.

\textbf{HEARTBEAT\_INSTRUMENTATION\_PLAN.md (350+ lines).}
Specifies a minimal-overhead circular buffer (115\,ns per tick, 0.01\%
overhead). Defines delta metric tracking, integration points in
\texttt{vm.c} (\texttt{vm\_init}, \texttt{vm\_cleanup},
\texttt{vm\_heartbeat\_run\_cycle}), and a CSV export function. Addresses
storage and sampling for up to 75~million tick rows.

\textbf{PHASE\_2\_EXECUTIVE\_SUMMARY.md (267 lines).}
Provides stakeholder-level clarity on the paradigm shift, the seven-metric
feedback web, the three measurement scales (tick, run, configuration), and
the golden-configuration selection criteria (five weighted dimensions).

\subsection{Key Insights}

\begin{enumerate}
  \item The heartbeat is an active control variable, not a passive clock.
  \item Metrics must be measured simultaneously as a coupled system;
    isolation is not possible.
  \item 750 runs are required for correlation power; the experiment therefore
    runs six to eight hours.
  \item The optimal configuration is not the one with the highest mean
    performance but the one with the deepest attractor basin, fastest
    convergence, tightest load coupling, and lowest jitter.
\end{enumerate}

\subsection{Implementation Checklist at Session End}

Structs and data definitions were complete. Integration points in
\texttt{vm.c} were specified. The analysis pipeline (convergence detection,
coefficient of variation, correlation matrix, spectral decomposition,
composite stability scoring) was designed. Seven C-level code changes and
three R/Python analysis components remained unimplemented.
